\RequirePackage{fix-cm}
\documentclass[11pt,a4paper]{ltjsarticle}
\usepackage[margin=24mm]{geometry}
\usepackage{amsmath,amssymb,amsthm,mathtools,booktabs}
\usepackage[haranoaji]{luatexja-preset}
\usepackage[unicode,hidelinks]{hyperref}
\newtheorem{theorem}{定理}[section]
\newtheorem{proposition}[theorem]{命題}
\newtheorem{lemma}[theorem]{補題}
\newtheorem{corollary}[theorem]{系}
\theoremstyle{definition}\newtheorem{remark}[theorem]{注意}
\renewcommand{\proofname}{証明}
\newcommand{\OO}{\mathcal O}\newcommand{\cA}{\mathcal A}\newcommand{\cB}{\mathcal B}
\DeclareMathOperator{\qgr}{qgr}\DeclareMathOperator{\Hom}{Hom}\DeclareMathOperator{\Ext}{Ext}
\DeclareMathOperator{\Jac}{Jac}\DeclareMathOperator{\Spec}{Spec}
\title{三次曲面のAS heartとanticanonical heartの一致開集合\\\large 一頂点cornerの有限判定と有理数上の証明書}
\author{研究原稿：標識付きnoetherian AS領域での結果}
\date{2026年11月15日\quad corner-v7 / hearts-v8}
\begin{document}\maketitle
\begin{abstract}
三つの三頂点blockを持つcubicのrolled-up quiverについて，構造対象に対応する
頂点のcornerがAS heart全体を再構成する条件を有限次元商で判定する．
反標準周期のcut gradingと計算用の道長gradingを区別する．
商の一つの道長成分の消滅は全高次数消滅を強制し，対応する条件はpotential係数空間の開条件となる．
六つの有理点からpotentialを明示構成し，頂点を除いた商のHilbert関数
$(8,21,33,38,33,21,8,0,\ldots)$を有理数上で検証する．
従って7次の有限行列で定まる非空開集合が得られる．さらに18係数を四交比に還元し，九つの小行列式と二つの行列式の非零性がcorner有限次元性の必要十分条件になることを証明する．この開集合を標識付き可換三次曲面の六点配置空間と明示的に同定する．さらに道長七の商が一様に零となることと，全 Hilbert 関数が $8+21t+33t^2+38t^3+33t^4+21t^5+8t^6$ に固定されることを証明する．開集合上に階数 $162$ の有限局所自由な代数族が得られる．
この開集合のnoetherian AS領域では，標準helixの構造対象から作った反標準環の
$qgr$とAS heartは自然にabel同値となる．任意のpotentialのAS正則性，
VdB heart・Rains heartとの全域一致，別の導来同値の不存在は主張しない．
\end{abstract}

\section{対象・grading・既知入力}
基礎体は標数0の代数閉体とする．頂点を
$A_a,B_b,C_c$ ($0\le a,b,c\le2$)，矢を
\[
 x_{ab}:A_a\to B_b,\qquad y_{bc}:B_b\to C_c,\qquad z_{ca}:C_c\to A_a
\]
とする．道の積は右から左への合成で，potentialとJacobi環は
\[
 W(w)=\sum_{a,b,c}w_{abc}[z_{ca}y_{bc}x_{ab}],\qquad
 B(w)=\Jac(Q,W(w)),\qquad e=e_{A_0}
\]
である．このquiverとpotentialはAOU\cite[\S4]{AOU}のcubic型である．
$w$を動かす係数空間は$\mathbb A^{27}$．全ての矢の次数を1とする
\emph{道長grading}を$\ell$と書く．全関係は道長2である．
一方，\emph{cut grading}は
\[
 \deg x=\deg y=0,\qquad\deg z=1
\]
である．反標準環と$qgr$の比較には後者を用い，$B^{\mathrm{cut}}$と書く．
両gradingで環の非次数付き構造は同じである．各cut次数の道長には有限の上界があるので，
$B^{\mathrm{cut}}$も局所有限で，次数0部分は有限次元である．

AOU\cite[(3.19), (4.9), (4.13)]{AOU}の標準的なhelix同定を使うAS領域では，
\[
 \cA_{\mathrm{AS}}\simeq\qgr B^{\mathrm{cut}},\qquad
 O=\pi(eB),\qquad \tau=\mathbb S^{-1}[2]=(1)_{\mathrm{cut}}.
\]
ここではfoldingがこの同定を与える標準的なmarked AS対象を扱う．
AS次元3のprojectiveの飽和から，cut次数付き同型
\begin{equation}\label{eq:orbit}
 eB^{\mathrm{cut}}e\simeq R(O,\tau)
 =\bigoplus_{n\ge0}\Hom_{\cA_{\mathrm{AS}}}(O,\tau^nO)
\end{equation}
を得る．飽和の理由は，AS条件により単純対象からprojectiveへの
$\Ext^0,\Ext^1$が消え，有限長加群へ組成列で拡張できることにある．
$qgr$のHomを定める有限長の切断を取ってもHomは変わらない．
右加群のorbit積は$ab=\tau^{\deg b}(a)\circ b$とする．

\begin{remark}
道長gradedの$\qgr B$を，そのまま一つのAS heartと同一視しない．
blockと道長の合同類による分解があるためである．
cornerの道長は3の倍数であり，反標準次数はその3分の1である．
cut gradingでfoldingすればこの曖昧さはない．
一般の係数でAS正則性やSerre同定を新たに証明したのではなく，
上記の標準AS領域に限定して適用する．
\end{remark}

\section{有限次元corner商と開性}
\begin{lemma}\label{lem:finite}
$B$は有限quiverの道長斉次関係による商で，$e$は頂点の和とする．
$C=B/BeB$に対して$C_N=0$なら，$C_n=0$ for $n\ge N$である．
このとき$Be$は右$eBe$加群，$eB$は左$eBe$加群として有限生成である．
\end{lemma}
\begin{proof}
$C$は矢で生成されるので，$C_N=0$の後は乗法によって全成分が零になる．
従って道長$N$の道は，$e$を通る同じ長さの道の線形結合に書ける．
$Be$の長い道の終端側$N$個をこの形で置き換えると，
右側に$eBe$の元を取り出し，左側を長さ$\le N$の$Be$の道にできる．
$e$で切る位置が端にある場合も，短い側の長さを$N$まで許せばよい．
残る場合は長さの帰納法を使う．よって長さ$\le N$の有限個の道が右加群を生成する．
逆向きの道について同じ議論を行えば$eB$の左有限生成性を得る．
この証明はcut gradingにも斉次的である．
\end{proof}

\begin{proposition}\label{prop:corner}
$B^{\mathrm{cut}}$が右noetherianであるとする．$C=B/BeB$が有限次元なら，
自然な関手
\begin{equation}\label{eq:equiv}
 \qgr B^{\mathrm{cut}}\longrightarrow\qgr(eB^{\mathrm{cut}}e),\qquad M\longmapsto Me
\end{equation}
はabel同値である．逆に，このrestrictionが$qgr$で同値なら$C$は有限次元である．
\end{proposition}
\begin{proof}
$R=eBe$は右noetherianである．実際，$R$の右イデアル$I$に対し
$(IB)e=I$なので，右イデアルの増大列は$B$のnoetherian性で止まる．
補題\ref{lem:finite}より$Me$は有限生成$R$加群となる．
restrictionは完全で，有限次元加群を保つ．
左随伴$L(N)=N\otimes_R eB$は有限生成加群を有限生成加群に送り，
$eB$の左有限生成性から有限次元加群も保つ．
unitは$(N\otimes_R eB)e\simeq N$であり，counitの核と余核は$e$で消える．
核は$B$の右noetherian性で有限生成である．
両者は有限次元環$C$の有限生成加群なので有限次元である．
$L$が短完全列の左端に作る核も同じ理由で有限次元となるから，
$qgr$上で$L$は完全となり，unit/counitが同値を与える．
逆に$\pi C$はrestrictionで零に送られる．同値なら$\pi C=0$であり，
$C$は有限生成$B$加群なので有限次元である．
\end{proof}

$V_N$を$A_0$を除くquiverの道長$N$の道を基底とするベクトル空間とする．
全ての$u(\partial_aW)v$，$\ell(u)+\ell(v)=N-2$を列挙した係数行列を$M_N(w)$とする．
その成分は$w$の線形式であり，$C_N$はその行空間による$V_N$の商である．
\begin{theorem}\label{thm:open}
\[
 U_N=\{w:\operatorname{rank}M_N(w)=\dim V_N\},\qquad
 U=\bigcup_{N\ge2}U_N
\]
は係数空間の開集合であり，$U$は$\dim C(w)<\infty$と同値である．
さらに，ある一様な整数$N_0$が存在し$U=U_{N_0}$となる．
$U$上のnoetherian AS対象では
\[
 \cA_{\mathrm{AS}}\simeq\qgr R(O,\tau)
\]
であり，この同値は構造対象と反標準shiftを保つ．
従ってこの範囲では導来同値だけでなくheart自体が一致する．
\end{theorem}
\begin{proof}
最大minorの非消滅が$U_N$を定め，補題\ref{lem:finite}から
$U_N\subset U_{N+1}$となる．有限次元性はある$N$での消滅と同値．
係数空間$\mathbb A^{27}$はnoetherianなので開集合$U$は準コンパクトであり，
この昇鎖はある$N_0$で$U$を覆う．$N_0=7$とは主張しない．
命題\ref{prop:corner}と式\eqref{eq:orbit}を適用する．
$\pi(eB)$の像は$\pi R$であり，restrictionはcut shiftと可換する．
\end{proof}

\begin{remark}
これは\emph{指定したcorner restriction}の一致開集合である．
$U$の外で別のabel同値や導来同値が存在しないとは結論しない．
全potentialがAS正則，noetherianであるとも，AS locus全体が有限型代数stackであるとも
本定理では仮定・主張しない．射の定義域を同じ有限係数chartで固定した場合の結論である．
\end{remark}


\section{係数トーラスの四変数化と無限次元cornerの障害}
\label{sec:torus}
$C(w)=B(w)/B(w)e_{A_0}B(w)$を考える．$A_0$に接する矢は$C(w)$で消えるので，
この商は$w_{abc}$ ($a=1,2$)の18係数だけに依存する．
それらが全て非零の範囲で
\[
q_{bc}=\frac{w_{2bc}}{w_{1bc}},\qquad
\lambda_{ij}=\frac{q_{ij}q_{00}}{q_{i0}q_{0j}}
\quad(1\le i,j\le2)
\]
と定める．
\begin{theorem}[四交比と九つの必要条件]\label{thm:torus}
\begin{enumerate}
\item 個々の矢の非零rescalingを除いた18係数の軌道は，
四つの$\lambda_{ij}\in k^\times$で完全に分類される．
\item $C(w)$が有限次元ならば18係数は全て非零であり，さらに
\begin{equation}\label{eq:minor}
q_{bc}q_{b'c'}-q_{bc'}q_{b'c}\ne0
\quad(b<b',\ c<c')
\end{equation}
の九条件が成立する．
\item $C(w)$が有限次元になるlocusは，
四次元係数トーラス上の非空の真の開集合である．
以前の可換六点witnessでは
\[
(\lambda_{ij})=
\begin{pmatrix}-1/4&1/16\\1/6&-1/4\end{pmatrix},\qquad C_7=0.
\]
\end{enumerate}
\end{theorem}
\begin{proof}
まず$y_{bc}$をrescaleして全$w_{1bc}$を1にできる．
残りの$x_{2b}/x_{1b}$および$z_{c2}/z_{c1}$のrescalingにより，
$q_{bc}$は任意の行因子・列因子を掛けられる．
$q_{00}=q_{b0}=q_{0c}=1$と正規化すれば
\[
Q=\begin{pmatrix}1&1&1\\1&\lambda_{11}&\lambda_{12}\\
1&\lambda_{21}&\lambda_{22}\end{pmatrix}
\]
だけが残る．逆に四交比はrescalingで不変なので，軌道を完全に記述する．

もし$w_{abc}=0$ ($a=1,2$)ならば，
$x_{ab},y_{bc},z_{ca}$だけを残し他の矢を零にすると
全てのJacobi関係が消える．各矢を$t$倍の行列単位に送り，
残る長さ3の閉路は非零な$t^3$倍の対角行列単位になる．
その全ての冪が非零なので$C(w)$は無限次元である．

全18係数が非零として式\eqref{eq:minor}の一つが零とする．
選んだ二行の添字集合$I$と二列の添字集合$J$で
$q_{bc}=r_bs_c$ ($b\in I,c\in J$)と分解できる．
行和・列和が零で全成分非零の
\[
T=\begin{pmatrix}1&-1\\-1&1\end{pmatrix}
\]
を$I\times J$に置く．$b\in I,c\in J$上で
\[
x_{1b}=1,\quad x_{2b}=-r_b^{-1},\quad
z_{c1}=1,\quad z_{c2}=s_c^{-1},\quad
y_{bc}=\frac{T_{bc}}{w_{1bc}}
\]
とし，残りの矢は零とする．Jacobi関係の左辺は
$T$の行和・列和，または
$w_{1bc}-w_{2bc}/(r_bs_c)$の倍数だからすべて零である．
各非零矢を$t$倍の行列単位に置き換えると，
支持に含まれる閉路が非零な$t^3$倍の対角行列単位となり，
$C(w)$は無限次元である．

有限次元性のlocusは前節の定理により開である．
四交比による正規化はこの条件を四次元トーラス上へ降ろす．
六点witnessの係数を代入すると表示した交比を得て，
既知の$C_7=0$より非空．
式\eqref{eq:minor}の零集合を含まないので真の開集合である．
\end{proof}

\begin{remark}[九つの明示的な禁止divisor]
四交比で式\eqref{eq:minor}は
\[
\begin{gathered}
\lambda_{11},\lambda_{12},\lambda_{21},\lambda_{22}\ne1,\\
\lambda_{11}\ne\lambda_{12},\ 
\lambda_{21}\ne\lambda_{22},\
\lambda_{11}\ne\lambda_{21},\
\lambda_{12}\ne\lambda_{22},\\
\lambda_{11}\lambda_{22}\ne\lambda_{12}\lambda_{21}
\end{gathered}
\]
に等しい．これら九つだけでは不十分であり，二つの行列式障害を加えた十一条件の必要十分性を定理\ref{thm:general-position-iff}で示す．AS locusとの交叉を調べることなしに
heartの非一致や導来非同値を結論しない．
\end{remark}
\section{行列式で検出する追加の無限次元障害}
\label{sec:ranktwo-obstructions}
前節の九つの二次小行列式条件だけでは，corner の有限次元性に
十分ではない．$w_{1bc}\ne0$ を仮定し，矢 $y_{bc}$ の rescaling により
$w_{1bc}=1$ として，$w_{2bc}=q_{bc}$ を記す．
\emph{成分ごとの逆数}からなる行列を
\[
 Q=(q_{bc})_{0\le b,c\le2},\qquad
 Q^{\langle-1\rangle}=(q_{bc}^{-1})_{0\le b,c\le2}
\]
と書く．後者は通常の逆行列 $Q^{-1}$ ではない．

\begin{proposition}[二つの行列式障害]\label{prop:two-determinants}
全ての $q_{bc}\in k^\times$ とする．もし
\[
 \det Q=0\quad\text{または}\quad
 \det Q^{\langle-1\rangle}=0,
\]
ならば $C(w)=B(w)/B(w)e_{A_0}B(w)$ は無限次元である．
従って $C(w)$ の有限次元性には，前節の九つの小行列式の
非零性に加えて，上の二つの行列式の非零性が必要である．
これら二つは九条件からは従わない．
\end{proposition}
\begin{proof}
まず $\det Q=0$ を仮定する．$Q$ を
\[
 Q=UV,\quad U\in M_{3\times2}(k),\quad V\in M_{2\times3}(k)
\]
と分解できる．$0\ne h\in\ker U^{\mathsf t}$ および
$0\ne k\in\ker V$ を選ぶ．それぞれが非零なので，
$h_bk_c\ne0$ となる $b,c$ が存在する．
$u_b$ を $U$ の第 $b$ 行，$v_c$ を $V$ の第 $c$ 列とする．

$k[t]$ 上の quiver 表現を
\[
 A_1=k[t]^2,\quad A_2=k[t],\quad B_b=C_c=k[t]
\]
とし，矢の写像を
\begin{equation}\label{eq:ranktwo-representation}
 x_{1b}=t u_b,\quad z_{c1}=-t v_c,\qquad
 x_{2b}=z_{c2}=t,\quad y_{bc}=t h_bk_c
\end{equation}
で定める．$x_{1b}$ は行ベクトル，$z_{c1}$ は列ベクトルである．
$y_{bc}$ に関する Jacobi 関係は
\[
 x_{1b}z_{c1}+q_{bc}x_{2b}z_{c2}
 =t^2(-u_bv_c+q_{bc})=0.
\]
$x_{1b},z_{c1}$ に関する関係はそれぞれ
$Vk=0$ と $h^{\mathsf t}U=0$ に帰着する．
$x_{2b},z_{c2}$ に関する関係は $Qk=0$ と
$h^{\mathsf t}Q=0$ に帰着する．
従って \eqref{eq:ranktwo-representation} は全ての Jacobi 関係を満たす．
しかも $A_2$ に基点を持つ三角閉路
\[
 z_{c2}y_{bc}x_{2b}:A_2\longrightarrow A_2
\]
の像は $t^3h_bk_c\ne0$ であり，その任意の冪の像も非零である．
よって $C(w)$ に非零元が無限に存在し，無限次元となる．
ここで $k[t]$ は表示の標的であり，有限次元の quiver 表現を
仮定していない．

次に $\det Q^{\langle-1\rangle}=0$ とする．
頂点 $A_1,A_2$ および対応する $x,z$ を交換すると，
残る二系列の potential 係数は $(q_{bc},1)$ となる．
各 $y_{bc}$ を非零スカラーで rescale して第一系列を1に戻せば，
新しい比行列は $Q^{\langle-1\rangle}$ である．
この操作は頂点を入れ替えた graded quiver の同型なので，
前半を適用できる．
\end{proof}

\begin{corollary}[四交比座標での十一の必要条件]
\label{cor:eleven}
\[
 Q\sim
 \begin{pmatrix}1&1&1\\1&a&b\\1&c&d\end{pmatrix}
 \qquad(a,b,c,d\in k^\times)
\]
と規格化する．$C(w)$ が有限次元ならば，前節の九つの二次小行列式が
全て非零であり，さらに
\begin{equation}\label{eq:two-extra-polynomials}
 (a-1)(d-1)-(b-1)(c-1)\ne0,\qquad
 bc(1-a)(1-d)-ad(1-b)(1-c)\ne0
\end{equation}
が必要である．
\end{corollary}
\begin{proof}
最初の式は $\det Q$ である．第二の式は
$abcd\det Q^{\langle-1\rangle}$ に等しい．
前節の定理と命題\ref{prop:two-determinants}を適用する．
\end{proof}

\begin{remark}[九つの小行列式だけでは不十分]
\label{ex:ranktwo}
\[
 Q_0=\begin{pmatrix}3&4&5\\5&6&7\\7&8&9\end{pmatrix},\quad
 Q_1=2520Q_0^{\langle-1\rangle}
 =\begin{pmatrix}840&630&504\\504&420&360\\360&315&280\end{pmatrix}.
\]
両方とも9つの $2\times2$ 小行列式は全て非零である．
一方，
\[
 \det Q_0=0,\quad
 \det Q_0^{\langle-1\rangle}=\frac1{198450},\qquad
 \det Q_1=80640,\quad
 \det Q_1^{\langle-1\rangle}=0.
\]
従って，両方とも $C$ は無限次元である．これらの厳密な有理数演算と
式\eqref{eq:ranktwo-representation}の全 Jacobi 関係の検査は
\texttt{computations/verify\_cubic\_corner\_rank2.py}で再現できる．
\end{remark}

\begin{remark}[有限体での探索と限界]
十一の非零条件の十分性は，次節の定理\ref{thm:general-position-iff}で六点配置との同定によって証明する．
規格化した四交比トーラスを $\mathbb F_7$ 上で全探索すると，
十一の条件を満たす80点すべてで $C_7=0$ になる．
これは有限体の有限個の点に関する独立検算である．標数0の十分性は次節の幾何学的証明に依存し，この有限体計算に依存しない．
同じスクリプトのオプション \texttt{--scan-f7} で再現できる．
\end{remark}

\section{四交比 corner の完全判定と六点配置による十分性}
\label{sec:six-point-classification}
本節では，前節までに証明した十一の必要条件が十分でもあることを示す．
その鍵は，四交比トーラスを\emph{標識付き可換三次曲面の六点配置空間}と同定することである．
以下，基礎体は標数0の代数閉体とし，十八の係数$w_{1bc},w_{2bc}$はすべて非零とする．

\begin{theorem}[corner有限性の必要十分条件]
\label{thm:general-position-iff}
$B(w)=\operatorname{Jac}(\widetilde Q,W(w))$，$e=e_{A_0}$，
$C(w)=B(w)/B(w)eB(w)$と置く．
$Q=(q_{bc})$，$q_{bc}=w_{2bc}/w_{1bc}$とすると，次は同値である．
\begin{enumerate}
\item $\dim_kC(w)<\infty$．
\item $Q$の九つの$2\times2$小行列式がすべて非零であり，
$\det Q\ne0$かつ$\det Q^{\langle-1\rangle}\ne0$である．
\item $P=(Q^{\langle-1\rangle})^{\mathsf t}$の三つの列$p_1,p_2,p_3$と，
座標点$p_4=(1:0:0),p_5=(0:1:0),p_6=(0:0:1)$が，
$\mathbb P^2$の一般位置の六点をなす．
\end{enumerate}
従って，四交比トーラスにおける有限corner locusは，
\emph{正確に}前節の十一のdivisorの補集合である．
この開集合は，六点の順序を保つ標識付き可換三次曲面のモジュライ
$M_{\mathrm{cubic}}$と自然に同型になる．
\end{theorem}
\begin{proof}
(i)$\Rightarrow$(ii)は定理\ref{thm:torus}と
命題\ref{prop:two-determinants}による．
(ii)$\Leftrightarrow$(iii)を示す．$P$の各成分は$q_{bc}^{-1}\ne0$であり，
その$2\times2$小行列式は，添字と非零係数を除けば
$Q$の対応する小行列式と等しい．さらに
$\det P=\det Q^{\langle-1\rangle}$である．従って(ii)のうち
九小行列式と後者の行列式は，順に「可動点二つと座標点一つが非共線」
「可動点三つが非共線」を表す．非零成分は「可動点一つと座標点二つが非共線」を表す．
座標点三つはもとより非共線である．

残る条件は六点が一つの二次曲線上にないことである．
座標三点を通るconicの方程式は
$\alpha X_1X_2+\beta X_2X_3+\gamma X_3X_1=0$
であり，可動点$p_i=(P_{0i}:P_{1i}:P_{2i})$への評価行列は
\[
 M_{i\bullet}=(P_{0i}P_{1i},P_{1i}P_{2i},P_{2i}P_{0i})
 =(P_{0i}P_{1i}P_{2i})
 (P_{2i}^{-1},P_{0i}^{-1},P_{1i}^{-1}).
\]
すべての因子が非零なので$\det M\ne0$は
$\det\bigl((P^{\langle-1\rangle})^{\mathsf t}\bigr)=\det Q\ne0$
と同値である．これで(ii)$\Leftrightarrow$(iii)を得る．

最後に(iii)$\Rightarrow$(i)を証明する．
$S=\operatorname{Bl}_{p_1,\ldots,p_6}\mathbb P^2$は滑らかな可換三次曲面で，
$H=-K_S$はampleである．AOU\cite[\S\S4,8，特に(4.1)，(4.7)，(4.9)，(8.4)--(8.7)]{AOU}
の九線束$E_v$とそのrolled-up helix代数を取る．
その幾何学的potentialの二つの関係系列は，基底を適切に規格化すると
\begin{equation}\label{eq:geometric-relations}
 w^{\mathrm{geo}}_{1bc}=P_{cb},\qquad
 w^{\mathrm{geo}}_{2bc}=1\qquad(0\le b,c\le2)
\end{equation}
となる．これは次の初等的な式からも直接確認できる．
$u_b=(P_{0b}P_{1b},P_{1b}P_{2b},P_{2b}P_{0b})$，
$L_{b,c}(X)=\det[p_{b+1},p_{c+4},X]$とし，
$D_b(X)=(u_{b+1}\times u_{b+2})\cdot
(X_1X_2,X_2X_3,X_3X_1)$（添字を$3$で読む）とする．
$D_b$は$p_{b+1}$以外の五点を通るconicである．
また$M_b(X)=\det[p_{b+2},p_{b+3},X]$と書くと，
行列式の展開により全ての$c$について
\begin{equation}\label{eq:six-point-identities}
 \sum_{b=0}^2P_{cb}D_b(X)L_{b,c}(X)=0,\qquad
 \sum_{b=0}^2M_b(X)L_{b,c}(X)=0
\end{equation}
が成立する．各関係空間はAOU(4.2)により一次元であるため，
\eqref{eq:six-point-identities}が\eqref{eq:geometric-relations}を与える．
これらの多項式恒等式は独立の整数係数による厳密検算コードにも収めた．

\eqref{eq:geometric-relations}より，その二系列の比は
$w^{\mathrm{geo}}_{2bc}/w^{\mathrm{geo}}_{1bc}
=P_{cb}^{-1}=q_{bc}$である．
元の二系列と幾何学的二系列は各矢$y_{bc}$の非零rescalingで一致する．
残りの$w_{0bc}$は頂点$A_0$を除いた商$C$に影響しないから，
\[
 C(w)\simeq C\bigl(w^{\mathrm{geo}}\bigr)
\]
が道長次数付き代数として成り立つ．

幾何学的代数については，AOU(4.9)により
\begin{equation}\label{eq:geometric-rolled-orbit}
 B\bigl(w^{\mathrm{geo}}\bigr)
 \simeq\bigoplus_{v,v'}\bigoplus_{n\ge0}
 \operatorname{Hom}_S(E_v,E_{v'}(nH))
\end{equation}
がcut次数付き代数として成立し，$e=e_{0,0}$に対応する線束は
$E_{0,0}=\mathcal O_S$である．
任意の$v$について$m\gg0$を固定すれば
$E_v^\vee(mH)$は大域生成される．評価写像の核を$K_v$として
\[
 0\longrightarrow K_v\longrightarrow
 H^0(S,E_v^\vee(mH))\otimes\mathcal O_S
 \longrightarrow E_v^\vee(mH)\longrightarrow0
\]
を取る．これを$E_{v'}((n-m)H)$でtensorし，
Serre消滅$H^1(S,K_v\otimes E_{v'}((n-m)H))=0$（$n\gg0$）を使うと，
乗法写像
\begin{equation}\label{eq:serre-factorization}
 H^0(S,E_v^\vee(mH))\otimes
 H^0(S,E_{v'}((n-m)H))
 \twoheadrightarrow
 \operatorname{Hom}_S(E_v,E_{v'}(nH))
\end{equation}
が全射となる．左辺の積は必ず
$E_v\to\mathcal O_S(mH)\to E_{v'}(nH)$という，
構造頂点$e$を通る経路に対応する．頂点は有限個なので全ての$v,v'$に共通の
十分大きい$N$を取れば，
$B_n\subset BeB$（$n\ge N$）が従う．
cut次数零は有限次元で，各斉次成分も有限次元であるから，
$\dim_kC\bigl(w^{\mathrm{geo}}\bigr)<\infty$となる．
これで(i)が証明された．

最後に，標識付き六点のmoduliで座標三点$p_4,p_5,p_6$を固定すると，
残る自由度は$P$の各列のスカラーと座標三点を保つ対角$\mathrm{PGL}_3$で割ることである．
この商はちょうど$Q$の行・列rescaling商であり，
$P=(Q^{\langle-1\rangle})^{\mathsf t}$およびその逆は
全成分非零のトーラス上で正則だから，主張した開集合の同型が得られる．
\end{proof}

\begin{corollary}[AS heart と anticanonical heart の一致領域]
\label{cor:complete-as-ac}
十八係数非零の係数チャートで，標準AOU型のfolding，
反Serre作用のcut shiftとの同定，必要なnoetherian性を満たすAS対象に限る．
このとき自然なcorner restrictionが
\[
 \mathcal A_{\mathrm{AS}}\xrightarrow{\sim}
 \mathcal A_{\mathrm{ac}}=\qgr R(\mathcal O,\tau)
\]
という構造対象・反Serre作用と両立するabel同値であることと，
定理\ref{thm:general-position-iff}の十一条件は同値である．
\end{corollary}
\begin{proof}
定理\ref{thm:general-position-iff}と命題\ref{prop:corner}，
式\eqref{eq:orbit}を適用する．
\end{proof}

\begin{remark}
定理\ref{thm:general-position-iff}は，\emph{元の27係数potential全体が}
可換三次曲面から生じると主張するものではない．
$w_{0bc}$は任意でよく，幾何学的潜在係数と一致させる必要はない．
同定しているのは$A_0$除去後のcorner商である．
また，十一条件に違反した場合に任意の別のabel同値・導来同値が存在しないという主張でもない．
Rains heart，VdB heartとの比較は依然として独立の課題である．
\end{remark}

\section{一般位置における一様な消滅次数は七である}
\label{sec:uniform-c7}
本節は定理\ref{thm:general-position-iff}の有限次元性を定量化する．
次数とは特に指定しない限り道長次数を指し，cut次数と区別する．
有限体における既存の計算は証明の入力にしない．

\begin{theorem}[最良の一様消滅次数]\label{thm:uniform-c7}
標数零の代数閉体上，十八の係数$w_{1bc},w_{2bc}$が非零であり，
$Q=(w_{2bc}/w_{1bc})$とする．次は同値である：
\begin{enumerate}
\item $Q$の九つの$2\times2$小行列式，$\det Q$，
$\det Q^{\langle-1\rangle}$がすべて非零．
\item $C(w)=B(w)/B(w)e_{A_0}B(w)$が有限次元．
\item 道長次数七の成分$C(w)_7$が零．
\end{enumerate}
これらの条件の下で$C(w)_n=0$ $(n\ge7)$，
一方$\dim_kC(w)_6\ge8$である．従って四交比トーラス上
\[
 U=U_7=\{\Delta_{11}\ne0\},
\]
かつ七は最良の一様消滅次数である．$\Delta_{11}$は
十一の禁止divisorの定義多項式の積を表す．
\end{theorem}
\begin{proof}
(iii)$\Rightarrow$(ii)は$C$が矢と頂点で生成されることから従い，
(ii)$\Rightarrow$(i)は既存の必要条件である．
定理\ref{thm:general-position-iff}で六点を復元し，
$S=\operatorname{Bl}_{p_1,\ldots,p_6}\mathbb P^2$と置く．
$h$を直線の引き戻し，$e_i$を六つの例外曲線とし，
\[
\begin{aligned}
H&=3h-\sum_{i=1}^6e_i=-K_S, &
D_1&=-2h+\sum_{i=1}^6e_i,\\
D_2&=-h+e_1+e_2+e_3, &
q_j&=2h-\sum_{r\ne j}e_r \quad(4\le j\le6).
\end{aligned}
\]
$q_j$は五点を通るconicの真の変換で，
$q_j^2=-1$，$H\cdot q_j=1$．三blockの線束は
$\mathcal O(D_a)$ $(a=1,2)$，$\mathcal O(e_i)$ $(1\le i\le3)$，
$\mathcal O(h-e_j)$ $(4\le j\le6)$と，
構造頂点$A_0$の$\mathcal O_S$である．
AOU\cite[(4.9)]{AOU}の幾何的rolled-up helix代数を使う．
道長七の始点・終点は三種類（二十一組）だけであり，
対応する切断は
\[
\begin{array}{c|c|c}
\text{始点と終点}&\text{組数}&\text{対応するdivisor}\\\hline
 A_a\to B_i\ (a=1,2;\ 1\le i\le3)&6&2H+e_i-D_a\\
 B_i\to C_j\ (1\le i\le3;\ 4\le j\le6)&9&2H+h-e_i-e_j\\
 C_j\to A_a\ (4\le j\le6;\ a=1,2)&6&3H+D_a-h+e_j .
\end{array}
\]
それぞれの全切断が途中で$\mathcal O_S(mH)$を通ることを示す．

まず，四つの線束
\[
\begin{aligned}
L^-_1&=H-D_1=5h-2\sum e_r,&
L^-_2&=H-D_2=4h-2(e_1+e_2+e_3)-(e_4+e_5+e_6),\\
L^+_1&=H+D_1=h,&
L^+_2&=H+D_2=2h-e_4-e_5-e_6
\end{aligned}
\]
は，いずれも六本の互いに素な$(-1)$曲線を縮約する
$S\to\mathbb P^2$の直線束である．
実際，$L^-_1$は$q_1,\ldots,q_6$，
$L^-_2$は$h-e_1-e_2,h-e_1-e_3,h-e_2-e_3,q_4,q_5,q_6$，
$L^+_1$は$e_1,\ldots,e_6$，
$L^+_2$は$e_1,e_2,e_3,h-e_4-e_5,h-e_4-e_6,h-e_5-e_6$
を縮約する．従って$h^0(L_a^\pm)=3$．
$L^-_a\cdot e_i=2$ $(i\le3)$，
$L^+_a\cdot q_j=2$ $(j\ge4)$であり，
これらの$(-1)$曲線は縮約されず，平面の滑らかなconicへ
同型に写る．したがって
\begin{equation}\label{eq:c7-conic-restriction}
H^0(L^-_a)\xrightarrow{\sim}H^0(e_i,\mathcal O_{e_i}(2)),
\qquad
H^0(L^+_a)\xrightarrow{\sim}H^0(q_j,\mathcal O_{q_j}(2))
\end{equation}
は同型である．

\paragraph{$A_a\to B_i$（六組）．}
$N=2H+e_i-D_a$について
\[
0\to\mathcal O(2H-D_a)\to\mathcal O(N)
\to\mathcal O_{e_i}(2)\to0 .
\]
$3H-D_a$はampleなので
$H^1(2H-D_a)=0$（Kodaira消滅）．
左端の切断は$A_a\to\mathcal O(2H)\to B_i(2H)$を経由する．
$H^1(H)=0$から$H^0(H+e_i)\to H^0(e_i,\mathcal O_{e_i})$
は全射である．式\eqref{eq:c7-conic-restriction}を併用すると，
$A_a\to\mathcal O(H)\to B_i(2H)$を経由する切断は
$H^0(e_i,\mathcal O_{e_i}(2))$全体に制限される．
従って二種類の積が$H^0(N)$を張る．

\paragraph{$B_i\to C_j$（九組）．}
$L=h-e_j$はbase-point-free pencilである．
$3H-e_i-h+e_j$はnefかつbigである．
実際，二十七本の$(-1)$曲線との交点はすべて非負で，
自己交点は八である（del Pezzoの錐定理を用いる）．
Kawamata--Viehweg消滅で$H^1(2H-e_i-h+e_j)=0$．
pencilの評価列を$\mathcal O(2H-e_i)$でtensorして
\[
H^0(2H-e_i)\otimes H^0(h-e_j)
\twoheadrightarrow H^0(2H+h-e_i-e_j)
\]
を得る．右辺は全て$B_i\to\mathcal O(2H)\to C_j(2H)$を経由する．

\paragraph{$C_j\to A_a$（六組）．}
$N=3H+D_a-h+e_j=2H+D_a+q_j$とする．
$3H+D_a$はampleだから$H^1(2H+D_a)=0$で，
\[
0\to\mathcal O(2H+D_a)\to\mathcal O(N)
\to\mathcal O_{q_j}(2)\to0
\]
は切断について右端まで全射である．左端は
$q_j$の定義切断を掛けた
$C_j\to\mathcal O(H)\to A_a(3H)$の積で張られる．
$H^1(H)=0$と$(H+q_j)\cdot q_j=0$から
$H^0(H+q_j)\to H^0(q_j,\mathcal O_{q_j})$が全射である．
式\eqref{eq:c7-conic-restriction}より
$H^0(H+D_a)\to H^0(q_j,\mathcal O_{q_j}(2))$は同型．
従って$C_j\to\mathcal O(2H)\to A_a(3H)$の積は
右端の全切断を張り，先の積と合わせて$H^0(N)$を生成する．

以上より，全二十一組で道長七の全元が$BeB$に属し，
$C_7=0$である．

\paragraph{道長六が残ること．}
八頂点$A_1,A_2,B_1,B_2,B_3,C_4,C_5,C_6$の
各対角の道長六の全空間は
$H^0(2H)$で，Riemann--RochとKodaira消滅により次元十．
$A_a$では$H\cdot D_a=0$と$H$のamplenessから
$H^0(\pm D_a)=0$である．
従って構造頂点を経由する積は$\mathcal O(H)$経由に限られ，
\[
H^0(H-D_a)\otimes H^0(H+D_a)\longrightarrow H^0(2H)
\]
の像の次元は高々$3\cdot3=9$である．
$B_i$では$B_i\to\mathcal O$はなく，
$\mathcal O(H)$経由の積を$e_i$へ制限すると
$H^0(H-e_i)$の高々二次元の像となる．
$\mathcal O(2H)$経由の積は$e_i$上で零．
一方，$H^1(2H-e_i)=0$から$H^0(2H)$は
$H^0(e_i,\mathcal O_{e_i}(2))$（三次元）へ全射．
$C_j$では$C_j\to\mathcal O$はなく，
$\mathcal O(H)$経由の積は$q_j$上で零，
$\mathcal O(2H)$経由の積を$q_j$へ制限すると
$H^0(h-e_j)$の高々二次元の像となる．
$H^1(2H-q_j)=H^1(H+h-e_j)=0$なので
$H^0(2H)$は$H^0(q_j,\mathcal O_{q_j}(2))$（三次元）へ全射．
従って八頂点の対角商は各々非零で，
$\dim C_6\ge8$．
\end{proof}

\begin{corollary}[道長七の関係行列の根基]\label{cor:c7-radical}
$\Lambda=k[a^{\pm1},b^{\pm1},c^{\pm1},d^{\pm1}]$とし，
規格化したpotentialの道長七のJacobi関係行列を$M_7$，
その列数（六千八百四）の最大小行列式のイデアルを$I_7$とする．
このとき
\[
\sqrt{I_7}=\sqrt{(\Delta_{11})}\qquad\text{in }\Lambda .
\]
\end{corollary}
\begin{proof}
定理\ref{thm:uniform-c7}より，
$M_7$の列full rank locusは$\{\Delta_{11}\ne0\}$に等しい．
従って$V(I_7)=V(\Delta_{11})$であり，
Laurent多項式環上のHilbert Nullstellensatzによって根基が一致する．
これは根基の等式であり，非被約scheme構造や個々のminorの
因数分解を述べていない．
\end{proof}

\section{一般位置六点における corner の Hilbert 関数の一様性}
\label{sec:hilbert-constant}
\begin{theorem}[corner の完全な Hilbert 関数]\label{thm:hilbert-constant}
標数零の代数閉体上，十八の係数 $w_{1bc},w_{2bc}$ が非零で，
定理\ref{thm:general-position-iff}の十一の非零条件を満たすとする．
$C=B/BeB$ に道長 grading を入れると，六点配置には依存せず
\[
 \boxed{H_C(t)=8+21t+33t^2+38t^3+33t^4+21t^5+8t^6,\qquad
       \dim_k C=162.}
\]
特に定理\ref{thm:uniform-c7}の $\dim C_6\ge8$ は等号であり，
$C_0,\ldots,C_6$ の次元跳びは一般位置 locus に存在しない．
この結論は $B$ 全体が AS 正則であることを仮定しない．
\end{theorem}

\begin{proof}
まず定理\ref{thm:general-position-iff}により，
$C$ を一般位置の六点を吹き上げた滑らかな可換三次曲面 $S$ の
rolled-up helix 代数の同じ idempotent 商として計算できる．
既存の記号を保ち，$H=-K_S$，
\[
\begin{gathered}
 D_1=-2h+\sum_{r=1}^6e_r,\qquad D_2=-h+e_1+e_2+e_3,\\
 A_0=\mathcal O_S,\quad A_a=\mathcal O_S(D_a)\ (a=1,2),\quad
 B_i=\mathcal O_S(e_i)\ (1\le i\le3),\\
 C_j=\mathcal O_S(G_j),\qquad G_j=h-e_j\ (4\le j\le6),\qquad
 q_j=H-G_j=2h-\sum_{r\ne j}e_r.
\end{gathered}
\]
ここで $q_j$ は五点を通る conic の真の変換である $(-1)$ 曲線の類を表す．
$A,B,C$ の三 block にレベル $0,1,2$ を与えると，
道長 $n$ の始点 $s$，終点 $t$ の全 Hom は，
AOU\cite[(4.1), (4.9)]{AOU}の同定により
\begin{equation}\label{eq:hilbert-block}
 B_n(s,t)=H^0\!\left(S,\mathcal O_S\left(
 \left\lfloor\frac{\ell(s)+n}{3}\right\rfloor H+E_t-E_s\right)\right)
 \quad(\ell(t)\equiv\ell(s)+n\pmod3),
\end{equation}
ただし $E_{A_a}=D_a$，$E_{B_i}=e_i$，$E_{C_j}=G_j$ である．
AOU の helix の acyclicity と曲面の Riemann--Roch
$\chi(\mathcal O(D))=1+D(D+H)/2$ を使うと，
$e=A_0$ を除く八頂点の全 Hom の各次数の合計は次の表の中央左列になる．
\[
\begin{array}{c|rrrrrrrr}
\text{道長 }n&0&1&2&3&4&5&6&7\\\hline
\dim B_n\text{（八頂点間）}&8&21&42&74&105&147&206&252\\
\dim (BeB)_n\text{（同上）}&0&0&9&36&72&126&198&252\\
\dim C_n&8&21&33&38&33&21&8&0
\end{array}
\]
以下では中央右列，すなわち構造頂点を\emph{途中で経由する}
積の像の次元を各次数について証明する．
多重経由があるときは，すべての経由位置に対応する像の和を取る．

必要な切断数を列挙する（$a=1,2$，$i=1,2,3$，$j=4,5,6$）：
\begin{equation}\label{eq:hilbert-basic-h0}
\begin{gathered}
 h^0(H)=4,\quad h^0(2H)=10,\quad h^0(H\pm D_a)=3,\\
 h^0(H-e_i)=2,\quad h^0(H+e_i)=5,\quad h^0(2H-e_i)=7,\\
 h^0(G_j)=2,\quad h^0(q_j)=1,\quad
 h^0(H+G_j)=7,\quad h^0(H+q_j)=5.
\end{gathered}
\end{equation}
$H-e_i$ と $G_j$ は base-point-free pencil である．
$H\pm D_a$ は定理\ref{thm:uniform-c7}で使った平面への blow-down 線束である．

\paragraph{道長二と三．}
道長二では $C_j\to B_i$ の九組のみが $e$ を経由し，
その像は $H^0(q_j)H^0(e_i)$ の一次元である．合計九次元．
道長三では $A_a\to A_b$ に経由射はなく，
$B_i\to B_k$ と $C_j\to C_k$ の九組ずつで，
それぞれ $H^0(H-e_i)H^0(e_k)$ および
$H^0(q_j)H^0(G_k)$ が二次元の像を与える．合計三十六次元．

\paragraph{道長四．}
全二十一組の全 Hom はそれぞれ五次元である．
$A_a\to B_i$（六組）では $H^0(H-D_a)H^0(e_i)$ の三次元が経由する．
$B_i\to C_j$（九組）では
$H^0(H-e_i)H^0(G_j)$ が経由する．後者の pencil trick の核は
$H^0(q_j-e_i)$ であり，これは零である．実際 $q_j\cdot e_i=1$ と
$q_j^2=-1$ から，そのような effective divisor は $q_j$ を成分として含むが，
差は $-e_i$ となって effective でない．従って積の像は四次元．
$C_j\to A_a$（六組）では $H^0(q_j)H^0(H+D_a)$ の三次元が経由する．
この三種類を加えると $6\cdot3+9\cdot4+6\cdot3=72$ である．

\paragraph{道長五．}
全二十一組の全 Hom はそれぞれ七次元である．
$A_a\to C_j$（六組）では
$H^0(H-D_a)H^0(G_j)$ の像を取る．$G_j$ の pencil trick の核は
$H^0(H-D_a-G_j)$ であり，零である．
$a=1$ なら $H-D_1-G_j=2q_j-e_j$ で，
$q_j$ との交点数は $-2$ だから $q_j$ が二重の固定成分となり，
残りが $-e_j$ となって矛盾する．$a=2$ なら
$H-D_2-G_j=3h-2(e_1+e_2+e_3)-\sum_{r=4,r\ne j}^6e_r$．
これは互いに素な三本の $(-1)$ 曲線
$h-e_1-e_2$，$h-e_1-e_3$，$h-e_2-e_3$ とそれぞれ交点数 $-1$ であり，
三本を引くと $-\sum_{r=4,r\ne j}^6 e_r$ となり effective でない．
よって積の像は $3\cdot2=6$ 次元．
$B_i\to A_a$（六組）では
$H^0(H-e_i)H^0(H+D_a)$ の像を取る．$H-e_i$ の pencil trick の核は
$H^0(D_a+e_i)=0$（$h$ との交点が負）なので六次元．
$C_j\to B_i$（九組）は $e$ を二通りの位置で経由し，
$H^0(H+q_j+e_i)$ の中で
\[
 q_j\,H^0(H+e_i)+e_i\,H^0(H+q_j)
\]
を得る．ここで曲線の記号はその定義切断を掛ける作用も表す．
二つの部分空間は共に五次元であり，$q_j$ と $e_i$ に共通成分がないので
交叉は $q_je_i\,H^0(H)$ の四次元である．よって和は六次元．
従って経由射の合計は $(6+6+9)\cdot6=126$ 次元である．

\paragraph{道長六：$A_a\to A_b$（四組）．}
この場合の経由射は一種類で，次の積の像である：
\begin{equation}\label{eq:two-blowdowns-injection}
 H^0(H-D_a)\otimes H^0(H+D_b)
 \longrightarrow H^0(2H+D_b-D_a).
\end{equation}
この積写像が全ての $a,b=1,2$ に対して単射であることを示す．
$L=H-D_a$，$M=H+D_b$ とする．いずれも $S\to\mathbb P^2$ の
blow-down を定め，切断空間は三次元である．
$L$ が縮約する六本の $(-1)$ 曲線のうち，$M\cdot E=2$ となる曲線は
$a=b$ で六本，$a\ne b$ で三本（後者は $q_4,q_5,q_6$）ある．
$L$ により縮約されるそれら三本の像は $\mathbb P^2$ で非共線である．
もし共線なら，その直線の真の変換は $H$ と交点零の有効曲線となり，
$H$ の ampleness に矛盾する．
一方，$M\cdot E=2$ なら，$M$ による blow-down は $E$ を
滑らかな平面 conic に同型に写す．したがって
\[
 H^0(M)\xrightarrow{\sim}H^0(E,M|_E)\simeq H^0(\mathbb P^1,\mathcal O(2))
\]
である．
仮に \eqref{eq:two-blowdowns-injection} の核に $\sum_{u=1}^3s_u\otimes t_u$
（$s_u$ は $H^0(L)$ の基底）があれば，三本の曲線への制限から
$\sum_{u=1}^3s_u(p_E)t_u=0$ が $H^0(M)$ で成り立つ．
三つの $p_E$ が非共線なので，三つの評価汎関数は独立であり，
全ての $t_u$ が零である．これで積写像の単射性が従う．
従って全ての四組で経由射の像は九次元である．

\paragraph{道長六：$B_i\to B_k$（九組）．}
$H^0(H-e_i)H^0(H+e_k)$ の像は，
$H-e_i$ の pencil trick と $h^0(e_i+e_k)=1$ より九次元である．
$i\ne k$ では全 Hom が九次元なのでこれで終わる．
$i=k$ のときは二番目の経由位置から $e_i\,H^0(2H-e_i)$ も得られるが，
$H^0(H-e_i)H^0(H)\twoheadrightarrow H^0(2H-e_i)$ が成立する．
実際，pencil trick の核は $H^0(e_i)$ の一次元で，
$2\cdot4-1=7=h^0(2H-e_i)$ だからである．
従って二番目の積は最初の像に含まれ，九次元のままである．

\paragraph{道長六：$C_j\to C_k$（九組）．}
$H^0(H+q_j)H^0(G_k)$ の像は，$G_k$ の pencil trick と
$h^0(q_j+q_k)=1$ より九次元である．
$j\ne k$ では全 Hom が九次元なのでこれで終わる．
$j=k$ の場合，もう一つの経由位置から $q_j\,H^0(H+G_j)$ も得るが，
$H^0(H)H^0(G_j)\twoheadrightarrow H^0(H+G_j)$ が成立する．
実際，pencil trick の核は $H^0(q_j)$ の一次元で，
$4\cdot2-1=7=h^0(H+G_j)$ である．
従って後者の積も先の九次元の像に含まれる．

以上より道長六の経由射の次元は
$(4+9+9)\cdot9=198$ である．
表中の全 Hom の次元 $206$ との差は八であり，$\dim C_6=8$ が従う．
道長七以降の消滅は定理\ref{thm:uniform-c7}による．
表を合計して Hilbert 関数と $\dim C=162$ を得る．
\end{proof}

\begin{corollary}[162次元 algebra の有限平坦族]
\label{cor:flat-corner-162}
標数零の代数閉体 $k$ 上の四交比トーラスにおける十一条件の補開集合
$U=\{\Delta_{11}\ne0\}$ 上，規格化した universal quiver potential の
$e_{A_0}$-corner 商を $\mathcal C_U$ とする．このとき
$\mathcal C_U$ は階数 $162$ の有限局所自由な
$\mathcal O_U$-代数で，その道長斉次部分は順に
\[
 (8,21,33,38,33,21,8)
\]
の階数を持ち，道長七以上は零である．任意の $U$ 上の基底変更と可換する．
\end{corollary}
\begin{proof}
各道長 $n$ の成分は有限個の道を基底とする自明ベクトル束を，
全ての $u(\partial W)v$ からなる有限個の関係によって割った
有限表示 $\mathcal O_U$ 加群である．
定理\ref{thm:hilbert-constant}より $0\le n\le6$ では
全幾何学的ファイバーの次元が一定であるから，
対応する関係行列は一定階数を持ち，商は局所自由で基底変更と可換する．
$n=7$ では全ファイバーが零なので，Nakayamaの補題で
$\mathcal C_{U,7}=0$．頂点と矢による生成性から高次数も零になる．
有限個の斉次部分を足し合わせて主張を得る．
\end{proof}

\begin{remark}[結果の対象範囲]\label{rem:hilbert-status}
証明は四交比の十一条件による一般位置六点の再構成と，
可換三次曲面の幾何学的 rolled-up helix の Hom 同定に基づく．
$w_{0bc}$ は $C$ に影響しないので任意でよい．
しかし $27$ 係数の任意の potential が AS-regular/noetherian であること，
Rains heart や VdB heart と同一であること，導来圏の全域での同値は
ここからは結論しない．
\end{remark}

\section{根壁と anticanonical heart の古典的障害}
\label{sec:root-walls}
本節は第\ref{sec:six-point-classification}節の四交比トーラスの
境界を古典的な六点 blow-up で解釈する．任意の非可換 potential の
AS 性，あるいは Rains heart・VdB heart の境界特殊化を仮定しない．
\begin{proposition}[十一の divisor の根幾何]
\label{prop:root-dictionary}
$P=(Q^{\langle-1\rangle})^{\mathsf t}$ の列を $p_1,p_2,p_3$，
座標点を $p_4,p_5,p_6$ とする．十八係数がすべて非零のとき，
十一の禁止条件は次の有効根に対応する．
\begin{enumerate}
\item $Q$ の二行 $b,b'$ と二列 $c,c'$ による $2\times2$ minor は，
$\{0,1,2\}\setminus\{c,c'\}=\{j\}$ と書けば
$p_{b+1},p_{b'+1},p_{4+j}$ の共線性であり，
根は $h-e_{b+1}-e_{b'+1}-e_{4+j}$．全部で九つ．
\item $\det Q^{\langle-1\rangle}=0$ は $p_1,p_2,p_3$ の共線性で，
根は $h-e_1-e_2-e_3$．
\item $\det Q=0$ は六点が一つの conic 上にある条件で，
根は $2h-\sum_{i=1}^6e_i$．
\end{enumerate}
各根 $\alpha$ は $\alpha^2=-2$ と $H\cdot\alpha=0$
（$H=3h-\sum e_i$）を満たす．各 divisor の一般点の
真の変換は滑らかな有理 $(-2)$ 曲線 $E$ である．
\end{proposition}
\begin{proof}
二行二列の minor の消滅は，$P$ の二列と補集合の座標点を
並べた $3\times3$ 行列式の消滅と同値である．
$\det P=\det Q^{\langle-1\rangle}$ であり，
座標三点を通る conic の三可動点での評価行列は，
非零因子を除いて $Q$ に一致する（定理\ref{thm:general-position-iff}）．
交点数は $h^2=1,e_i^2=-1$ から従う．
一般点では余分な共線性・重複点・conic の退化はない．
\end{proof}

\begin{theorem}[一つの根壁では導来圏も異なる]
\label{thm:classical-root-wall}
十一因子の一般点に対応する六点 blow-up を
$S=\operatorname{Bl}_{p_1,\ldots,p_6}\mathbb P^2$，
対応する $(-2)$ 曲線を $i:E\hookrightarrow S$ とする．
$H=-K_S$ と置く．
\begin{enumerate}
\item $S$ は滑らかな弱 del Pezzo 曲面で $H^2=3$，
$H$ は nef かつ big だが ample ではない．反標準射
\[
f:S\longrightarrow X=\operatorname{Proj}\bigoplus_{n\ge0}H^0(S,nH)
\]
は $E$ を収縮し，$X$ は $A_1$ 特異点を持つ正規三次曲面である．
\item $F=i_*\mathcal O_E(-1)$ は非零な $2$-spherical object であり，
\[
F\otimes\mathcal O_S(H)\simeq F,\qquad
\operatorname{RHom}_S(\mathcal O_S(nH),F)=0\quad(n\in\mathbb Z),
\qquad \mathbf Rf_*F=0.
\]
従って $(\mathcal O_S,(-)\otimes\mathcal O_S(H))$ は
$\operatorname{coh}S$ において Artin--Zhang ample でない．
\item $R_S=\bigoplus_{n\ge0}H^0(S,nH)$ と置くと
$\operatorname{qgr}R_S\simeq\operatorname{coh}X$ だが，
\[
\boxed{D^b(\operatorname{coh}S)\not\simeq
       D^b(\operatorname{qgr}R_S).}
\]
従ってこの古典的境界は，導来同値を前提とする
単なる HRS tilt の例ではない．
\end{enumerate}
\end{theorem}
\begin{proof}
十一因子の一般点では六点は相異なり，四点共線などの
追加退化はない．六点 blow-up に対する弱 del Pezzo の判定と，
反標準射が $(-2)$ 曲線を rational double point に収縮する定理を使う
\cite[Introduction, \S4]{BlancLamy}
\cite[Theorem 8.3.2]{Dolgachev}．
他の有効根のない一般点なので特異点は $A_1$ 型である．
$H\cdot E=0$ から $H$ は ample ではない．

$E\simeq\mathbb P^1$ かつ $E^2=-2$ より
$N_{E/S}\simeq\mathcal O_E(-2)$ である．
局所 Ext 層 $\mathcal Ext_S^0(F,F)=\mathcal O_E$，
$\mathcal Ext_S^1(F,F)=N_{E/S}$ とその局所大域スペクトル列から
\[
\operatorname{Ext}^j_S(F,F)=
\begin{cases}k&j=0,2,\\0&j\ne0,2.\end{cases}
\]
$K_S|_E\simeq\mathcal O_E$ なので Serre 関手は
$F\mapsto F[2]$ となり，$F$ は $2$-spherical．
同じ理由で $H|_E\simeq\mathcal O_E$ だから
\[
\operatorname{RHom}_S(\mathcal O_S(nH),F)
\simeq\mathbf R\Gamma(E,\mathcal O_E(-1))=0
\]
がすべての $n$ で成り立つ．
非零対象 $F$ が全反標準 twist から見えないので
AZ の生成条件は失敗する．
$f(E)$ は一点で $H^0(E,\mathcal O_E(-1))
=H^1(E,\mathcal O_E(-1))=0$ より $\mathbf Rf_*F=0$．

$X$ は $A_1$ 特異点を持つ正規三次曲面であり，
その homogeneous coordinate ring の $\operatorname{qgr}$ は
$\operatorname{coh}X$ である．滑らかな $S$ では
有界 coherent 複体の任意の二対象間の高次 Ext は消滅する．
ところが $X$ の特異点 $x$ の非正則局所環では
剰余体 $k(x)$ の最小自由解消が無限に続き，
$\operatorname{Ext}_X^j(k(x),k(x))\ne0$ となる $j$ が
無限に存在する．この性質は三角圏の同値と両立しない．
\end{proof}

\begin{corollary}[任意次数の弱 del Pezzo に対する障害]
\label{cor:all-degree-weak-dp}
標数零の代数閉体上の任意次数 $d=K_S^2\ge1$ の滑らかな
弱 del Pezzo 曲面 $S$ に $(-2)$ 曲線 $E$ があるとする．
$\tau=(-)\otimes\mathcal O_S(-K_S)$ に対し
$(\mathcal O_S,\tau)$ は AZ ample ではなく，
\[
D^b(\operatorname{coh}S)\not\simeq
D^b\!\left(\operatorname{qgr}
 \bigoplus_{n\ge0}H^0(S,-nK_S)\right).
\]
\end{corollary}
\begin{proof}
任意次数の弱 del Pezzo の反標準モデルは $(-2)$ 曲線を
Du Val 特異点へ収縮する\cite[Theorem 8.3.2]{Dolgachev}．
同じ $F=i_*\mathcal O_E(-1)$ の全反標準 twist からの
導来 Hom が零となること，および反標準モデルの特異点で
剰余体が非有界の Ext を持つことを使う．
\end{proof}

\begin{proposition}[一つの根壁の有理数上の明示例]
\label{prop:explicit-single-root-wall}
\[
Q_0=\begin{pmatrix}1&1&1\\1&2&2\\1&3&4\end{pmatrix}
\]
では行 $(0,1)$・列 $(1,2)$ の minor だけが零であり，
$\det Q_0=1$，$\det Q_0^{\langle-1\rangle}=1/24$．
六点は
\[
(1:1:1),(2:1:1),(12:4:3),(1:0:0),(0:1:0),(0:0:1).
\]
共線三点は $p_1,p_2,p_4$ だけで，直線は $y=z$．
$|-K_S|$ の基底を
\[
\begin{aligned}
 A&=\tfrac13x^2(y-z)-xy(y-z),\\
 B&=\tfrac7{12}x^2(y-z)-x(y^2-z^2),\\
 C&=\tfrac16x^2(y+2z)-\tfrac32xy^2+y^2z,\\
 D&=\tfrac14x^2(y+z)-\tfrac32xy^2+yz^2
\end{aligned}
\]
とすると反標準像の三次方程式は
\[
\begin{aligned}
0={}&23A^3-22A^2B+44A^2C-42A^2D+8AB^2-16ABC\\
&-24AC^2+24AD^2+8B^2C+24BC^2-24BCD .
\end{aligned}
\]
曲線 $y=z$ の像は $[A:B:C:D]=[0:0:1:1]$．
この点で勾配は零であり，$C=1$ の局所座標
$(A,B,D-1)$ の Hessian 行列式は $-2304\ne0$．
従って実際に通常二重点である．
\end{proposition}

\begin{remark}[非可換への適用範囲]\label{rem:noncommutative-wall-scope}
本節は古典的 blow-up の coherent heart とその
anticanonical heart の比較である．十一因子上の一般の非可換
potential が AS heart を定めるか，Rains heart・VdB heart の
特殊化が $\operatorname{coh}S$ と同定できるかは未証明．
特に VdB の一般位置での導来同値をこの根壁まで
無条件に延長してはならない．本例では
$\mathbf Rf_*$ 自身が非零対象を零に送るので，
単なる振幅一の同値に伴う HRS tilt とは根本的に異なる．
\end{remark}

\section{非空性の明示的な証明書}
\subsection{六点からpotentialを復元する}
次の点を$\mathbb P^2$に取る：
\[
 p_1=(1:0:0),\ p_2=(0:1:0),\ p_3=(0:0:1),\quad
 p_4=(1:1:1),\ p_5=(1:2:3),\ p_6=(1:4:2).
\]
全三点行列式は非零であり，二次単項式
$(x^2,xy,xz,y^2,yz,z^2)$の六点評価行列式は$-12$である．
従って三点は共線でなく，六点はconic上にない．そのblow-up $S$は滑らかなcubicである．

AOUのfoundationの三blockを
\[
 \begin{gathered}
 (\OO,\OO(-2h+\textstyle\sum e_i),\OO(-h+e_1+e_2+e_3)),\\
 (\OO(e_1),\OO(e_2),\OO(e_3)),\qquad
 (\OO(h-e_4),\OO(h-e_5),\OO(h-e_6))
 \end{gathered}
\]
とする\cite[(4.1)]{AOU}．
$\ell_{ij}$を$p_i,p_j$を通る直線，$q_i$を$p_i$以外の五点を通るconicとする．
係数は原始整数列で最初の非零成分を正として規格化する．
$A_a\to B_b\to C_c$の三つの合成切断は，$i=b+1,j=c+4$と置くと
\[
 \ell_{ij},\qquad q_i\ell_{ij},\qquad
 \ell_{\{1,2,3\}\setminus\{i\}}\ell_{ij}
\]
で表される（左から$a=0,1,2$）．
各$a,c$に対する三切断の唯一の関係を同じ原始規格化で取り，$w_{abc}$とする．
これはAOUのrelation-extension potentialの直接構成\cite[(4.7)]{AOU}である．
得られる行列$W_a=(w_{abc})_{b,c}$は
\begin{equation}\label{eq:witness}
 W_0=\begin{pmatrix}1&1&1\\-1&-2&-2\\1&3&1\end{pmatrix},\quad
 W_1=\begin{pmatrix}2&1&1\\-1&2&-4\\-1&-3&1\end{pmatrix},\quad
 W_2=\begin{pmatrix}1&1&2\\-1&-1&-1\\1&1&1\end{pmatrix}.
\end{equation}
各関係は多項式係数として厳密に検証する．Euler行列だけからの推測ではない．
この可換点でのAS性はAOUのhelixの既知結果を用い，有限rankから推論しない．
noetherian性は，rolled algebraが反標準section環上の有限加群であることから従う．
実際，九線束の各Hom束の十分高次のsectionは有限生成であり，低次数も有限次元である．

\subsection{商の全関係行列}
$A_0$を除くと，頂点8，矢21，二次関係21が残る．
残る全関係を全てのprefixとsuffixで掛け，始点・終点別に行列を作る．
\begin{proposition}[厳密な有限計算]\label{prop:certificate}
式\eqref{eq:witness}のpotentialに対し，道長gradingで
\[
 H_{B/BeB}(t)=8+21t+33t^2+38t^3+33t^4+21t^5+8t^6.
\]
特に$w_*\in U_7$であり，$U_7$は非空である．
\end{proposition}
\begin{proof}
再現コード\texttt{computations/verify\_cubic\_corner.py}は，点から式\eqref{eq:witness}を復元し，
全Jacobi関係を微分から構成する．$C_n$のpresentation行列は全ての
$u(\partial W)v$を列挙しており，二次生成イデアルの次数$n$部分を尽くす．
有理数上の分数演算による行消去で次を得る：
\begin{center}
\begin{tabular}{c|rrrrrrrr}
道長$n$&0&1&2&3&4&5&6&7\\\hline
道の数&8&21&54&144&378&972&2592&6804\\
$\dim C_n$&8&21&33&38&33&21&8&0
\end{tabular}
\end{center}
同じ行列を素数$1000003$で還元した別の消去でも同じrankを得る．
7次は21個の始終点blockの全てで列full rankであり，計6804列を消去する．
finite-fieldでのfull rankだけでも整数minorの非零を証明するが，低次数の正確な次元には
有理数計算を用いている．浮動小数点は用いない．
$C$は道長1で生成されるので$C_7=0$から全高次数が零となり，表示が全Hilbert級数となる．
総次元は$162$である．証明書の行列生成と検算は計算機に依存するが，
新しいparameterについてのAS性をこの検算で置き換えてはいない．
\end{proof}

以上により，標識と構造頂点を固定したnoetherian AS領域で，AS heartとanticanonical heartが
一致する明示的な非空開集合を得た．係数は全て非零なのでpotentialの係数torusとも交わる．
矢の非零rescalingは$e$を保つため，有限次元性と$U_N$はその作用で不変である．
従って標識を保つ商にも同じ開条件が降りる．

\section{HRS比較を進める符号判定}
前稿の振幅1比較を$F:D^b(\cA)\simeq D^b(\cB)$と書き，
$F(\cA)\subset D^{[0,1]}(\cB)$，$F^{-1}(\cB)\subset D^{[-1,0]}(\cA)$とする．
$(T,U)=(\ker H^1F,\ker H^0F)$はtorsion pairである．
その理由を記す．$M\in\cA$について$B_i=H^iFM$とし，truncation triangleを
$G=F^{-1}$で移すと$GB_0\to M\to GB_1[-1]$となる．
振幅条件から$H^{-1}GB_0=H^0GB_1=0$であり，
$0\to GB_0\to M\to H^{-1}GB_1\to0$を得る．
端の二項はそれぞれ$T,U$に属し，$\Hom(T,U)=0$は
$F(T)\subset\cB$，$F(U)\subset\cB[-1]$から従う．
従って$G(\cB)=\langle U[1],T\rangle$である．
\begin{proposition}\label{prop:sign}
$F(e)=e'$，$F\tau=\tau'F$で，両作用が対応するheartを保つとする．
$(e',\tau')$がAZ-ampleで高次Serre消滅を満たすと仮定する．
$0\ne M\in\cA$が$\tau M\simeq M$かつ単純なら，$\chi(e,M)\ne0$であり，
\[
 M\in T\ \Longleftrightarrow\ \chi(e,M)>0,\qquad
 M\in U\ \Longleftrightarrow\ \chi(e,M)<0.
\]
$M\in U$なら$\Ext^i(e,M)$は$i=1$だけで非零となり，
$\dim\Ext^1(e,M)=-\chi(e,M)$である．
\end{proposition}
\begin{proof}
単純性とtorsion分解により$M$は$T$または$U$に属する．
$N=H^0FM$または$N=H^1FM$は非零で$\tau'$固定である．
高次消滅と固定性から$\Ext^{>0}(e',N)=0$．
もし$\Hom(e',N)=0$なら全twistからのHomが零となりAZ生成性に反する．
従って$h=\dim\Hom(e',N)>0$．$FM=N$ならEuler値は$h$，
$FM=N[-1]$なら$-h$であり，後者のRHomは次数1に集中する．
\end{proof}
この符号判定は，VdBの各段階でpointednessと反標準作用の適合性を検証した後に適用する．
共鳴式だけから必要な単純対象の存在を作ったことにはならない．
一般の非単純対象について，Euler値が負というだけで対象全体が$U$に入るとはしない．

\section{結果の範囲と残る問題}
定理\ref{thm:open}は一般のcorner原理の証明と，cubicの有限表示での明示化である．
命題\ref{prop:certificate}の証明は有理数・有限体の厳密計算を含む．
外部独立査読と新規性の全面調査は未実施．

反標準環はAS側の指定対象$O$から作ったものなので，
AOU \S9のmutation後の先頭対象がVdB模型の$\pi F$と同一であるか，
および既知の特定の導来同値がこのcorner同値に一致するかは別問題である．
今回そこを仮定なしに解決したとはしない．
また$U_7$の外が全て非一致であるわけではなく，より高次の$U_N$に入ることがある．
Rains heartとVdB heart（後者はliteral blow-upに由来する），共鳴ごとのtorsion pair，特異境界は引き続き比較する．
数値的共鳴の稠密性から，これらの導来非同値の密度を推論してはいけない．

\begin{thebibliography}{9}
\bibitem{AOU} T.~Abdelgadir, S.~Okawa, K.~Ueda,
\emph{Compact moduli of marked noncommutative cubic surfaces},
arXiv:2404.00175v1 (2024), \S\S3--4, 9.
\url{https://arxiv.org/abs/2404.00175}.
\bibitem{OU} S.~Okawa, K.~Ueda,
\emph{Marked Artin--Schelter surfaces of del Pezzo types},
arXiv:2508.04986v1 (2025), Definitions 1.1, 5.2 and \S6.
\url{https://arxiv.org/abs/2508.04986}.
\bibitem{VdB} M.~Van den Bergh,
\emph{Blowing up of non-commutative smooth surfaces},
Mem. Amer. Math. Soc. 154 (2001), no.~734;
arXiv:math/9809116v1, Theorems 9.1.9, 11.1.3.
\url{https://arxiv.org/abs/math/9809116}.
\bibitem{BlancLamy} J.~Blanc, S.~Lamy,
\emph{Weak Fano threefolds obtained by blowing-up a space curve and construction of Sarkisov links},
Proc. London Math. Soc. (3) 105 (2012), 1047--1075, Introduction and \S4.
\url{https://doi.org/10.1112/plms/pds023}.
\bibitem{Dolgachev} I.~V.~Dolgachev,
\emph{Classical Algebraic Geometry: A Modern View},
Cambridge University Press (2012), Theorems 8.3.2 and 8.3.4.
\bibitem{Hearts} 本研究，\emph{非可換del Pezzo曲面のheart比較}，
\texttt{papers/nc\_del\_pezzo\_hearts.tex}，2026年11月15日．
HRS表示と一般corner判定の現行証明．本稿は有限生成性を道長判定から導く．
\end{thebibliography}
\end{document}
